\documentclass[12pt]{article}
\usepackage[utf8]{inputenc}
\usepackage[T1]{fontenc}
\usepackage[margin=1in]{geometry}
\usepackage{setspace}
\usepackage{xcolor}
\usepackage[colorlinks=true, urlcolor=blue!60!black, linkcolor=blue!60!black]{hyperref}
\setlength{\parindent}{0em}
\setlength{\parskip}{0.6em}
\renewcommand{\baselinestretch}{1.0}
\begin{document}
\begin{center}
{\LARGE \textit{Trump v. Anderson}, 601 U.S. 100 (2024)\par}
\vspace{0.5em}
{\large Case Brief --- Two-Page Edition\par}
\end{center}

\section{Memory Jogger}
States cannot enforce Section~3 of the Fourteenth Amendment against federal candidates, so Colorado could not remove Donald Trump from its presidential primary ballot (\href{https://www.supremecourt.gov/opinions/23pdf/601us1r06_m6hn.pdf}{601 U.S. 100}).

\section{Detailed Case Facts}
In September 2023, Colorado voters led by Norma Anderson petitioned under the state election code to keep Trump off the 2024 primary ballot, alleging that he organized and incited the January 6, 2021 breach of the Capitol and was disqualified under Section~3.

\section{Procedural History}
The trial court found that Trump ``engaged in insurrection'' but held that Section~3 does not cover the President. The Colorado Supreme Court, 4--3, ordered him off the ballot. The Supreme Court granted expedited review and reversed on March 4, 2024.

\section{Judicial Votes}
9--0 in the judgment: a per curiam joined by Roberts, Thomas, Alito, Gorsuch, and Kavanaugh; Barrett concurred in part; Sotomayor, Kagan, and Jackson concurred in the judgment.

\section{Holding}
States may disqualify persons holding or seeking state office but have no power to enforce Section~3 against federal officeholders and candidates. Five Justices added that enforcement against federal offices rests with Congress through Section~5 legislation.

\section{Analysis of Opinions}
The per curiam reasoned that the Fourteenth Amendment expanded federal power at the States' expense, so power over federal offices must be delegated, not reserved; state exclusions would burden Congress's power to lift disabilities by a two-thirds vote; history, including \textit{Griffin's Case}, shows Congress supplied enforcement procedures; and divergent state rulings would create a chaotic presidential patchwork. Barrett would not have decided whether legislation is the exclusive vehicle and urged the Court to ``turn the national temperature down.'' Sotomayor, Kagan, and Jackson agreed on the result but accused the majority of deciding much more than necessary and undercutting the two-thirds amnesty rule.

\section{Examples: Future Applications}
\textbf{Same-Side:} under Section~3, a State cannot remove a candidate for the Senate or the House. \textbf{Opposite-Side:} a State may bar a former sheriff from its governor's race on the same record. \textbf{Fence-Sitter:} a State excludes a twice-elected former President under the Twenty-Second Amendment.

\section{Critique}
Requiring legislation lets a bare majority nullify a provision that takes two-thirds of both Houses to lift; critics also fault the reliance on ``patchwork'' fears and a single circuit decision, and the Court left the President's status and the meaning of ``insurrection'' unresolved.

\section{Key Quotations}
\begin{quote}
``We conclude that States may disqualify persons holding or attempting to hold state office. But States have no power under the Constitution to enforce Section 3 with respect to federal offices, especially the Presidency.''
\end{quote}
\noindent --- Per Curiam, 601 U.S. at 110.

\end{document}
